% Research directions enabled by your multimodal Qur’an dataset

% I divide them into: (A) Speech-only, (B) Text-only / NLP, (C) Speech + Text (multimodal), and (D) Tools & Applications. For each I include: what to build, why your dataset helps, methods/approaches, evaluation ideas, and main challenges.

% A — Speech-only (audio-focused)
% A1 — Qur’anic ASR (verse- and word-level)

% What: Train automatic speech recognition systems specialized for Qur’anic recitation (word/verse transcription, segmentation).
% Why dataset helps: You have aligned word-level audio + text and many reciters → enables supervised ASR with variation across voices and recitation styles.
% Approach: End-to-end ASR (Conformer / Wav2Vec2 / Whisper-style fine-tuning) with language model boosting for Classical Arabic text; optionally adapt a grapheme-to-phoneme module for tajweed-influenced pronunciations. Use per-word audio for fine-grained segmentation models.
% Eval: Word error rate (WER) on verse and on token-level, boundary detection metrics (F1) for segmentation, per-reciter and cross-reciter WER.
% Challenges: Classical (Fusha) Qur’anic orthography differs from modern corpora; long melodic stretching (madd), pauses, and Quranic tajweed rules create unusual acoustic patterns that need modeling. See prior ASR overviews and datasets. 
% ResearchGate

% A2 — Automatic detection/classification of Tajweed rules (mispronunciation tagging)

% What: Detect occurrences of tajweed phenomena (e.g., madd, ghunnah, idgham, ikhfa, tafkheem/tarqeeq) or identify mistakes in learners’ recitations.
% Why dataset helps: Word-level audio aligned with canonical text allows supervised labels (you could annotate examples or derive weak labels from expert recitations). Existing smaller datasets (e.g., QDAT) target this task but are limited in size/diversity — your dataset can scale this. 
% ResearchGate
% +1

% Approach: Frame-level or segment-level classifiers (CNN/transformer on spectrograms; sequence labeling with CTC or token classification); combine acoustic features with phonetic/linguistic features (G2P outputs, rule-based signals). Use multi-task learning to jointly detect multiple tajweed phenomena.
% Eval: Precision/recall/F1 per rule; accuracy of error detection for beginner vs expert recitations; human expert validation.

% A3 — Reciter identification and voice profiling

% What: Given a verse or segment, identify the reciter (voice recognition) or cluster reciters by style/dialect.
% Why dataset helps: 32 reciters × consistent verse coverage → ideal for speaker ID, voice-style analysis, and voice-clustering research.
% Approach: Speaker embeddings (x-vectors, ECAPA-TDNN), contrastive learning across verses to capture style-invariant features; investigate style vs content disentanglement.
% Eval: Top-1/top-5 accuracy for reciter ID; clustering metrics (ARI, NMI) for style groupings; cross-reciter generalization tests.
% Challenges: Reciters sometimes change style or tempo; dataset does allow large-scale study though. Several recent papers use similar reciter datasets for speaker/classification tasks. 
% Natural Science Publishing

% A4 — Prosody & melodic analysis (intonation, madd duration modeling)

% What: Measure and model prosodic patterns, stretch durations, pause insertion patterns, and melodic contours of recitation styles.
% Why dataset helps: Multiple reciters reciting same verses lets you analyze prosody independent of lexical content.
% Approach: Extract F0, duration, intensity, align to tokens; use sequence models to predict expected madd lengths, or to synthesize recitation prosody for TTS.
% Eval: Correlation between predicted and real durations/F0; listening tests for synthesized prosody.

% B — Text-only (NLP / classical Arabic)
% B1 — Token-level morphological analysis & diacritization for Qur’anic Arabic

% What: Automatic diacritization (short vowels/harakat) and morphological parsing specialized for Classical Qur’anic Arabic.
% Why dataset helps: Word-level tokens with transliteration and aligned audio can supervise diacritization and pronunciation-to-text mapping.
% Approach: Sequence-to-sequence or transformer diacritization models; multimodal input (audio + orthography) to improve disambiguation.
% Eval: Diacritization error rate, morphological tagging accuracy.

% B2 — Semantic retrieval, question answering, and RAG (Qur’anic knowledge retrieval)

% What: Build retrieval-augmented systems to answer queries using Qur’anic text + translations + tafsir (if available).
% Why dataset helps: Fine-grained word-level translations + verse-level context support better chunking, retrieval and grounding for LLMs.
% Approach: Create embeddings (text-only or multimodal) per verse/word for vector search; fine-tune LLM or use RAG pipelines for QA grounded in the Qur’an.
% Eval: Retrieval metrics (recall@k), QA accuracy vs human-annotated answers, faithfulness metrics.

% B3 — Alignment of translation and transliteration to source tokens

% What: Create reliable mappings between Arabic tokens and English translations at word-level (not trivial: translation is often non-literal or many-to-one).
% Why dataset helps: You already have per-word English translations and transliterations that can be used to train alignment models or to produce better bilingual lexicons.
% Approach: Use attention-based alignment models or alignment via forced alignment (multimodal cues). Train word-level bilingual embeddings.
% Eval: Alignment precision/recall vs human alignments.

% C — Multimodal (speech + text together) — where your dataset excels
% C1 — Word-level forced-alignment and phoneme-level segmentation

% What: Accurate alignment between Arabic orthography and audio at word/phoneme level (time-stamps for each token).
% Why dataset helps: You have word-level audio files which can bootstrap or supervise alignment; many prior works lack this fine-grain mapping at scale.
% Approach: Use Kaldi / Montreal Forced Aligner or neural forced alignment with CTC-based models; train acoustic models using the canonical tokens as transcripts.
% Eval: Boundary error (seconds), percent tokens aligned correctly. This alignment is foundational for other downstream tasks (pronunciation modeling, TTS unit extraction).

% C2 — Pronunciation tutoring & automated feedback for learners

% What: Systems that listen to a learner recite, detect tajweed issues, and provide targeted, explainable feedback with audio examples.
% Why dataset helps: Canonical recitations from 32 expert reciters act as high-quality references; word-level audio enables precise corrective examples.
% Approach: Combine ASR + tajweed-classifier + sequence alignment to highlight mismatches; create explainable feedback modules with synthetic examples (show correct audio, visual articulatory cues).
% Eval: Expert-rated improvement in learner recitation; detection precision/recall for learner mistakes (user studies).

% C3 — Qur’an TTS that respects tajweed and recitation styles

% What: A TTS system that synthesizes Qur’anic recitation obeying tajweed rules and optionally mimics specific reciters’ melodic style.
% Why dataset helps: Parallel verse coverage for 32 reciters allows multi-speaker TTS training and style-transfer (voice cloning) while word-level audio helps unit selection or neural vocoder conditioning on pronunciation/pronunciation duration (madd).
% Approach: Modern neural TTS (FastSpeech2 + HiFi-GAN) with style embeddings, or fine-tune multi-speaker Tacotron/Glow-TTS; include tajweed-informed linguistic front-end (G2P + prosody markers). Consider neural prosody modeling conditioned on tajweed labels. Prior TTS/Quran attempts show feasibility but often with limited corpora. 
% Ijassat
% +1

% Eval: MOS for naturalness and tajweed correctness, expert evaluation for adherence to recitation norms.

% C4 — Cross-reciter style transfer and voice conversion with tajweed preservation

% What: Convert a recitation from one reciter to sound like another while preserving tajweed-correct pronunciation and prosody.
% Why dataset helps: Same textual content across many reciters gives parallel (or near-parallel) training data; word-level audio helps aligning units for conversion.
% Approach: Voice conversion methods (GAN/flow-based or neural-adaptive). Add constraints to preserve phonetic realization consistent with tajweed.
% Eval: Speaker similarity (speaker verification), intelligibility, and expert-rated tajweed adherence.

% C5 — Multimodal embeddings & downstream tasks (classification, retrieval)

% What: Learn joint audio-text embeddings for semantically-aware indexing (e.g., search by meaning + recitation style).
% Why dataset helps: Paired audio and text across many reciters allows multimodal representation learning that can capture both meaning and acoustic style.
% Approach: Contrastive multimodal pretraining (e.g., CLIP-style for audio-text), then use embeddings for retrieval, clustering, or transfer learning.
% Eval: Retrieval metrics, clustering coherence, downstream task gains (e.g., better ASR under low-resource conditions using multimodal pretraining).

% D — Applications, tooling & community resources
% D1 — Evaluation suites & benchmarks

% Create standardized test splits (per-surah, per-reciter) and tajweed-annotated test sets for community benchmarking (ASR, tajweed detection, TTS). Several recent works highlight the need for shared benchmarks. 
% arXiv

% D2 — Data augmentation and synthetic learner errors

% Generate controlled mispronunciations and synthetic learner recitations for robust training of error-detection and tutoring models. Papers synthesize errors for controlled evaluation in similar domains. 
% arXiv

% D3 — Tools for scholars: aligned corpora for linguistic / phonological research

% Offer easy-to-use APIs and visualizers (word audio players, aligned waveform + transcript viewers) to support humanities researchers and teachers.

% Suggested baseline experiments (practical starting points)

% ASR baseline — Fine-tune a pre-trained wav2vec2 model on your verse-level data; report WER per-reciter and averaged. Compare training on verse-level transcripts vs using word-level audio as additional fine-tuning data.

% Tajweed classifier baseline — Use QDAT-style labels (or create new labels) to train a CNN+LSTM classifier on spectrograms to identify a set of tajweed rules. Evaluate using F1. (QDAT and other papers provide a reference setup). 
% ResearchGate
% +1

% Forced-aligner baseline — Use Montreal Forced Aligner with an Arabic acoustic model trained on your word-audio to produce phoneme/time alignments; evaluate boundary error against manual checks.

% Multi-speaker TTS baseline — Train a multispeaker FastSpeech2 conditioned on reciter IDs and tajweed prosody tags; evaluate MOS + tajweed expert score.

% Evaluation metrics and human evaluation

% Automatic: WER/CER, precision/recall/F1 for tajweed classes, boundary error (ms), MOS (naturalness), speaker-ID accuracy, embedding retrieval metrics (recall@k).

% Human/Expert: Expert tajweed adherence score, MOS with domain experts, pedagogical usefulness in small user studies (does feedback help learners improve?). Prior work emphasizes expert validation for tajweed tasks. 
% ResearchGate
% +1

% Main challenges & gaps (research opportunities)

% Classical Qur’anic text vs Modern Arabic: Models pre-trained on modern corpora may mis-handle classical orthography, archaisms, and specific punctuation/pauses; domain-adaptation is needed.

% Tajweed-sensitive modeling: Tajweed rules create acoustic phenomena (madd, nasalization) that standard speech models don’t normally account for — specialized phonetic modeling and G2P with tajweed features are needed. 
% Ijassat

% Evaluation standardization: Lack of widely accepted benchmark splits and tajweed annotation conventions hinders cross-paper comparability — your dataset can host benchmark suites. 
% arXiv

% Cultural/ethical considerations: Recitation is a religious practice; synthetic recitations or voice cloning must be handled respectfully, with opt-in consent or clear purpose statements for the community.

\section{Future Work}

\textsc{Qur’an-MD} opens multiple avenues for future research at the intersection of speech processing, natural language processing, and digital Islamic studies. We outline three directions: (A) Speech-only tasks, (B) Multimodal approaches integrating speech and text, and (C) Tools and applications that can support both research and community engagement.  

\textbf{Speech-only Research}: A major direction lies in automatic speech recognition (ASR) for Qur’anic recitation, both at the verse and word level. The audio and verse-level transcripts aligned in the dataset with 30 reciters provide a unique testbed to develop end-to-end ASR models adapted to classical Arabic with the unique prosody and melodic contours in recitations. Evaluations can measure word error rate and segmentation accuracy across reciters, while challenges include modeling tajweed-driven prosody (madd, pauses, emphatics) and script-pronunciation mismatches.  

\textbf{Multimodal Speech-Text Research}: The full potential emerges when the modalities are combined. The dataset can drive progress in forced alignment and phoneme-level segmentation, producing high-quality timestamps for words and phonemes that benefit pronunciation modeling. It enables pronunciation tutoring systems tailored to Qur’anic recitation, which automatically detect tajweed issues and mispronounced words, providing corrective audio feedback and guidance on proper melodic and rhythmic patterns.

Another promising direction is Qur’anic TTS that respects tajweed rules and recitation styles. Leveraging parallel recordings from 32 reciters, the dataset enables the development of personalized TTS systems that provide corrective feedback for tajweed and Qur’anic melody in a user’s own voice, facilitating accurate imitation and learning. In addition, style transfer approaches can allow users to reproduce the prosody and melodic patterns of renowned reciters while preserving tajweed correctness, supporting both educational applications and expressive recitation synthesis. 

\textbf{Tools, Applications, and Community Resources:} \textsc{Qur’an-MD} also enables the development of advanced tools and applications that bridge research and community engagement. Multimodal embeddings of Qur’anic text and audio can support semantic retrieval of specific verses, style-aware search, and clustering of reciters, while facilitating the creation of robust tutoring systems for pronunciation and tajweed. Beyond machine learning applications, the dataset can be integrated into retrieval platforms, APIs, and interactive visualization tools for scholars, educators, and students, providing aligned verse- and word-level corpora that support linguistic, phonological, and prosodical research. Such resources can foster broader accessibility and engagement with the Qur'an and its studies.


% \textbf{Text-only / NLP Research}: On the text side, word-level annotations facilitate \emph{diacritization and morphological analysis} for classical Qur’anic Arabic, a long-standing challenge due to complex orthography. Incorporating aligned transliterations and pronunciations may further reduce diacritization errors. Another direction is \emph{semantic retrieval and question answering}, where embedded in aligned Arabic--English pairs can improve retrieval-augmented generation for Qur’anic queries. The dataset also enables research on \emph{alignment of translations and transliterations}, resulting in more accurate bilingual lexicons and improving cross-lingual applications in Islamic studies.  